\section{Scope and retrospective audit design}
\label{sec:design}

A stated privacy claim paired with its supporting experimental records forms
our basic unit of analysis. Scope covers four training campaigns and two
security-transfer arguments. Because we selected these episodes from an
earlier programme synthesis whose outcomes were already known, the resulting
archive traces how local measurements grew into sweeping claims while leaving
field-wide error rates unestimated. Component results remain credited to
their originating studies, and Table~\ref{tab:master-audit} summarises how
each recorded headline maps onto its audited disposition.

All historical runs finished before we specified this reconciliation.
Throughout the text, frozen exploratory protocols, completed pilot
executions, unexecuted confirmation proposals, and secondary consistency
checks remain explicitly separated. Our analysis recomputes stored loss
contrasts, probe advantages, and decision thresholds directly from archived
tables without retraining victim models or fitting stronger attackers.
Prospective validation requires an independent cohort. Where LLM-Attack
Suite defines observation contracts and control conventions
\cite{suite2026}, we apply its reporting structure to compare cross-campaign
claims under their original experimental budgets.

\paragraph{Evidence required for a claim.}
Auditing any candidate claim begins by logging the split protocol, victim
cohort, secret-key custody, remote host state, and observable tensor
channels. Target indexing and control bindings come next. Specifically, each
record pairs the protected token or property target with its exact tensor
index, training horizon, baseline guesser, acceptance gate, and independent
resampling unit. A qualification binding ties positive canary floors and
permutation-null bands to one concrete target, channel, defence family, and
probe implementation. Passing controls on one channel leaves unmeasured
channels unqualified.

Primary generated CSV and JSON outputs alongside active collection manifests
supersede informal narrative summaries. Whenever a subsequent technical
audit revises an earlier disposition, we report both the updated status and
every attached operational caveat regardless of which reading looks more
favourable. Every retained number carries explicit provenance.
Appendix~\ref{app:artifacts} binds each quantity to its hashed file path and
marks whether our scripts recomputed the statistic from raw seed records or
extracted it from a cited study.

\paragraph{Five outcomes for a specified claim.}
Individual endpoints within a single split configuration often merit
distinct evidentiary labels.
\begin{description}
  \item[Breach.] An admissible, calibrated probe exceeds its frozen
    acceptance threshold and establishes leakage for the specified observer,
    target, and finite protocol. Unprobed disclosures may still coexist
    alongside it.
  \item[Non-detection.] A calibrated attack stays within its acceptance rule
    under the tested architecture, dataset, and optimisation budget.
    Consistent with empirical differential-privacy auditing
    \cite{jagielski2020}, staying below an attack threshold leaves
    worst-case mutual information and adaptive recovery unbounded from above.
  \item[Undetermined.] Missing coverage or failed controls block privacy
    passes.
  \item[Invalid.] A verified implementation defect severs the recorded score
    from the claimed estimand, as occurs when next-token predictions are
    misread as current-token reconstructions. Properly relabelled, the
    underlying number may remain informative.
  \item[Out-of-scope.] Transferring a measurement to a stronger observer,
    altered key-custody regime, or untested protocol exceeds what the
    original experiment can support.
\end{description}

Utility loss and causal remote-learning gain are tracked alongside privacy
under each campaign's original conjunctive decision rule. Conjunctive gating
acts asymmetrically across endpoints. Within the 19-cell exploration, one
qualified breach or one utility failure suffices to reject a candidate even
when companion probes remain undetermined, whereas incomplete qualification
alone blocks any joint privacy pass. Failing an exploratory remote-gain
threshold similarly leaves population learning benefit unconfirmed on that
split.

\paragraph{Transferring a conclusion.}
Carrying an empirical privacy result into a new deployment setting requires
matching protocols, secret-key custody, observed tensors, target alignments,
and calibrated probe bindings. Wire confidentiality differs sharply from
host confidentiality. When an authorised server holds a shared orthogonal
key to restore activations in memory, passive wire non-detection says
nothing about a compromised host reading its own RAM \cite{permattack2025}.
Correlated token positions within fixed evaluation panels likewise add zero
independent training seeds and restrict how far finite exploratory
intervals generalise.

\begin{table}[t]
\centering\footnotesize
\setlength{\tabcolsep}{4pt}
\begin{tabularx}{\linewidth}{@{}>{\raggedright\arraybackslash}p{0.17\linewidth}>{\raggedright\arraybackslash}p{0.19\linewidth}>{\raggedright\arraybackslash}p{0.20\linewidth}>{\raggedright\arraybackslash}p{0.16\linewidth}>{\raggedright\arraybackslash}X@{}}
\toprule
Investigation & Evaluated view & Unaudited claim & Audited state & Retained quantitative basis \\
\midrule
Decoy gradients (\artifact{E1}) & Forward activations vs.\ joint gradient view & Passing forward check implies boundary privacy & Out-of-scope and Breach & Exact real/decoy row recovery on 4,096/4,096 frames per run (9 seeds), $+0.65$ to $+1.50\pp$ token gain, $0.01$-nat clip/noise utility cost \\
\addlinespace[3pt]
Boundary audit (\artifact{B1}) & Representations $(x_1,\ldots,x_{T-1})$ vs.\ $(x_2,\ldots,x_T)$ & Trained boundary resists input reconstruction & Invalid for current token & Shifted label scores prediction, 6/6 zero-return controls match or beat remote, answer match drops $34.79\%$ to $18.15\%$ \\
\addlinespace[3pt]
19-cell exploration (\artifact{R1}) & Inbound/outbound byte probes across 4 seeds & 19 rejections stem from 20/24 uncalibrated probes & Breach (13) and Undetermined (6) & 4/24 probes qualify, \RoneQualifiedBreaches\ breaches, \RoneUtilityFailures\ utility and \RoneGainFailures\ gain failures, \texttt{span-6} passes utility ($\Delta U=0.099971$) \\
\addlinespace[3pt]
Defence pilot (\artifact{P4}) & Structure and property probes across 3 paired seeds & $40\times$ excess implies broken gate or probe & Breach (4) and Undetermined (3) & \PfourUtilityFailures/\PfourCells\ fail utility, 7/7 pass gain, punctuation advantage $0.79$ to $0.89$ vs.\ $0.02$ gate \\
\addlinespace[3pt]
Wrapper and ratchet (\artifact{E8}, \artifact{ER}) & Passive wire capture under shared cloud key $Q$ & Zero-cost protection against compromised cloud & Non-detection (wire) and Out-of-scope (host) & $3.71\%$ pooled d1 wire score ($5.73\%$ max session), cloud holds $Q$, N=16 ratchet cuts tok/s by $63.16\%$ (20/24 identity) \\
\addlinespace[3pt]
Variational MI (\artifact{MI}) & DV critic readout $L_T$ on defended boundary & Near-zero MINE certifies small leak budget & Out-of-scope for upper ceiling & $L_T \le \operatorname{MI}(X,Z)$ bounds leakage from below only, $T=0$ yields $L_T=0$ when $Z=X$ leaks 8 bits ($m=256$) \\
\bottomrule
\end{tabularx}
\caption{Cross-campaign admissibility summary across all six audited
investigations. Each row pairs the evaluated observation surface and initial
reading with the primary audit finding and resulting evidentiary state from
Section~\ref{sec:design}. Hyperlinked record names point to the hashed
provenance entries in Appendix~\ref{app:artifacts}.}
\label{tab:master-audit}
\end{table}
